\section{Introduction}

We construct internal deductive systems from categorical axioms for
partial computation and prove essential incompleteness using inseparable
pairs. G\"odel's arithmetization starts with formulas and finite proofs
and represents their effective operations within arithmetic.
Di Paola and Heller proposed an intrinsic recursion theory supplying the
diagonalization and inseparability arguments needed for generalized
incompleteness \cite[pp.~594--596]{DPH}.

For a recursion category $\mathscr R$, let
$\mathcal D=\Tot(\Split(\mathscr R))$ be the category of total
morphisms after splitting the restriction idempotents, which represent
domains of partial computations. The deductive construction assumes that
$\mathcal D$ is regular, has disjoint pullback-stable finite coproducts,
and has the initial algebras for syntax, lists, and finite derivations
specified in Appendix~\ref{app:syntax-theories}. For an axiom subobject,
the image of the conclusion morphism, restricted to derivations of
sentences, is the generated sentence theory. When its inclusion has a
partial inverse, this image is represented by a restriction idempotent.
Its inverse image under sentence negation represents the refutable
sentences. Consistency means that the theorem and refutation domains are
disjoint. Completeness means that their join is the identity. We verify these existence and representation hypotheses in
the realizations below.

A consistent internal theory is essentially incomplete if every
consistent internal extension is incomplete. Two disjoint domains are
inseparable when no disjoint pair of restriction idempotents whose join
is the identity contains them, respectively. We prove two criteria.
The first uses a functor $F:\mathscr R\to\mathcal A$ into a split
restriction category that preserves restrictions, ranges, zero morphisms,
and binary restriction coproducts. Such a functor is called a range functor.
The functor extends to $\Split(\mathscr R)$ and interprets the theorem
and refutation domains in $\mathcal A$. Soundness, together with
disjointness of validity and its inverse image under negation, proves
these interpreted domains disjoint. Assume also that inverse images along
total morphisms in $\mathcal A$ preserve binary joins of restriction
idempotents. Under the hypotheses of
Theorem~\ref{thm:logical-comparison}, including a nonzero image of the
identity on the Turing object and a specified partial morphism in
$\mathcal A$ admitting a preimage under $F$, the comparison proves
incompleteness of both the interpreted theory and the internal theory.
The second criterion uses a total morphism mapping the two members of a
fixed inseparable pair into the theorem and refutation domains of a base
theory, respectively. This
reduction proves incompleteness of every consistent internal extension
(Corollary~\ref{cor:domain-incompleteness}).

Generation by finite deduction is left adjoint to the inclusion of
closed sentence theories. Logical translations preserving syntax and
inference rules induce Galois connections between theories ordered by
inclusion. Assume further that every monomorphism in $\mathcal D$ has a
partial inverse in $\Split(\mathscr R)$ and that the syntax and derivation
constructions exist for languages with nullary predicates indexed by any
object of $\mathcal D$. Theorem~\ref{thm:free-first-order-adjunction}
then constructs a
first-order theory whose nullary predicates are indexed by the object
on which a disjoint pair of domains is defined. Its
nonlogical axioms assert these predicates on one domain and their
negations on the other. If that theory is consistent, it is
essentially incomplete exactly when the pair is inseparable
(Corollary~\ref{cor:signed-first-order-incompleteness}).

For a sentence theory $T\hookrightarrow S$, the associated category
is the category of partial maps in the slice $\mathcal D/T$.
If the syntax and derivation constructions also exist in every slice
$\mathcal D/P$, for $P\in\mathcal D$, these are Turing categories,
and a theory inclusion induces a fully faithful functor between them
(Theorem~\ref{thm:theory-extension}). Suppose a fixed logical theorem is
represented by a total morphism $1\to S$ that factors through every
theory under consideration. A consistent theory in $\Split(\mathscr R)$
is essentially incomplete exactly when its corresponding internal theory
is essentially incomplete in the associated category. The internal theory
in the associated category defines a partial morphism that takes the value
true on its theorem domain and false on its refutation domain. This partial
morphism has no total extension in that category exactly when the theorem
and refutation domains of the original theory are inseparable in
$\Split(\mathscr R)$
(Theorem~\ref{thm:associated-incompleteness}). The second criterion
proves essential incompleteness and inseparability for the three
instances of this construction described below.

The arithmetic realization uses the syntactic recursion categories of
Montagna and Di Paola--Montagna \cite{Montagna,DPM}. For Peano arithmetic ($\PA$), the completion of their category
$S_{\PA}$ represents computably enumerable (c.e.) extensions of $\PA$
by internal theories. Let $M$ be the standard arithmetic model $\N$ or
an ultrapower of that model. Interpreting program graphs in $M$ gives a
range functor $F_M:S_{\PA}\to\Pfn$, where $\Pfn$ is the category of
sets and partial functions. A model of the represented axioms interprets the formulas
encoded by the sentence codes. Soundness and the comparison theorem
give incompleteness of both the internal theory and its interpretation
for each consistent c.e.~extension of PA
(Theorem~\ref{thm:arithmetic-internal-comparison}). We obtain recursively
inseparable theorem and refutation sets and an independent sentence
uniformly from an index for the axiom enumeration
(Corollary~\ref{cor:arithmetic} and
Theorem~\ref{thm:uniform-independent-sentences}). The associated-category
construction is applied separately to the first-order theory generated
by the arithmetic diagonal pair, whose nonlogical symbols are nullary
predicates (Corollary~\ref{cor:arithmetic-associated-categories}).

The second first-order realization uses a sequence
$a:\N\to[0,1]$ constructed by Miller \cite{Miller2004}.
We form a recursion category of partial numerical functions computed
by one oracle program with the same domain and values for every binary
representation of $a$. Its completion contains the internal deductive
system for a consistent theory of real closed fields expanded by
constants naming the coordinates of $a$. Every consistent extension is
incomplete if a single oracle program semidecides its axiom set using
any binary representation of $a$ as oracle
(Theorem~\ref{thm:miller-incompleteness}). This theory also supplies an
instance of the associated-category construction
(Corollary~\ref{cor:miller-associated-categories}).

The third realization uses cosheaves on a three-point space. The
recursion category constructed from them is not equivalent as a
restriction category to the category of partial recursive functions.
Its completion realizes propositional logic. Evaluation at the two
endpoints gives two consistent propositional theories, each of whose
consistent extensions in the same language with a c.e.~axiom set is
incomplete (Corollary~\ref{cor:cosheaf-essential-incompleteness}).
A conservative first-order extension of the internal propositional theory
supplies the third instance of
the associated-category construction
(Corollary~\ref{cor:fork-associated-theories}).

Section~\ref{sec:lifting} proves the computational comparison criteria,
and Section~\ref{sec:internal-theories} applies them to the sentence
theories constructed in Appendix~\ref{app:syntax-theories}.
Section~\ref{sec:theory-extensions} constructs the associated categories
and proves the extension and transfer theorems.
Sections~\ref{sec:syntactic} and~\ref{sec:models} establish the
syntactic realization and soundness. Sections~\ref{sec:arithmetic},
\ref{sec:miller}, and~\ref{sec:cosheaves} give the three applications.
The remaining appendices supply the cosheaf construction, the lifting
counterexample, and the faithful family of model interpretations.

\subsection{Related work}

Rosolini's $P$-categories generalized dominical categories
\cite{RosoliniThesis,Rosolini}; Heller adopted them and defined
prodominical categories \cite[Section~5]{Heller}.
We use the restriction categories of Cockett and Lack
\cite{CockettLack}. Binary restriction products give their equivalent
$P$-category formulation \cite[Proposition~4.5]{CockettLackIII}.
Heller's existence theorem \cite[Sections~8--9]{Heller} and
Lengyel's existence theorems using faithful iteration functors
\cite{Lengyel2004,Corrections} concern the construction of recursion
categories. Di Paola and Montagna's syntactic examples raised the
question of extending Heller's theorem beyond locally connected categories
\cite{DPM}.

Di Paola and Heller proved categorical recursive inseparability under
section and coproduct hypotheses
\cite[Theorem~8.15]{DPH}. Bernardi obtained categorical fixed-point
and inseparability results from precomplete equivalence relations,
including PA-provable equality of partial recursive functions among his
examples \cite[Section~3, pp.~125--128]{Bernardi}.
Montagna obtained effectively inseparable pairs in his syntactic
category \cite[Remark~3.1]{Montagna}.

Di Paola and Heller also discussed the role of the formalized L\"ob
theorem among the axioms of diagonalizable algebras and the modal
provability logic $\mathrm{GL}$,
and the obstacles to extending modal fixed-point theorems to numerals
and quantification \cite[pp.~595--596]{DPH}.

Cockett and Hofstra relate Turing categories to partial combinatory
algebras through idempotent splitting
\cite[Theorem~4.9]{CockettHofstra}. Cockett's later incompleteness
theorem assumes a predicate expressing program restriction order and
produces a restriction idempotent on the restriction terminal object
that is neither zero nor the identity
\cite[Section~7.5, Theorem~7.14]{CockettComputability}.
Its assumed predicate and conclusion differ from the construction of
generated sentence theories and the results for their consistent
extensions.

Bauer and Hanson use Miller's sequence in a parameterized realizability
topos with countable Dedekind reals \cite[Theorem~6.4]{BauerHanson}.
We use their binary representations to construct the recursion category
and prove essential incompleteness for the theory of real closed fields
expanded by constants.

